\documentclass[]{report}
% Language setting
% Replace `english' with e.g. `spanish' to change the document language
\usepackage[english]{babel}
% Set page size and margins
% Replace `letterpaper' with `a4paper' for UK/EU standard size
%\usepackage[letterpaper,top=2cm,bottom=2cm,left=3cm,right=3cm,marginparwidth=1.75cm]{geometry}
\usepackage[a4paper, margin=1in]{geometry}
% Useful packages
\setcounter{secnumdepth}{3}
\usepackage{amsmath}
\usepackage{xcolor}
\usepackage{indentfirst}
\usepackage{amssymb}
\usepackage{amsthm}
\usepackage{bbm}
\usepackage{pifont} 
\usepackage{booktabs}
\usepackage{bm}
\usepackage{comment}
\usepackage{graphicx}
\usepackage{subcaption} 
\usepackage{array}
\bibliographystyle{plain}
\usepackage{float}
\usepackage{lipsum} % For dummy text
\usepackage[colorlinks=true, allcolors=blue]{hyperref}
\theoremstyle{definition}
\newtheorem{theorem}{Theorem}
\newtheorem{definition}{Definition}
\newtheorem{example}{Example}
\newtheorem{remark}{Remark}
\newtheorem{lemma}{Lemma}
\newtheorem{proposition}{Proposition}
%opening
% \title{Random effects for high-dimensional correlated multinomial regression}

\begin{document}
	
\thispagestyle{empty}

\begin{itemize}
    \item Expectation-Maximization algorithm is a widely used iterative algorithm for finding the (local) maxima of the likelihood with incomplete data. It iteratively estimates the unknown parameters of the model by increasing the expected likelihood of the complete data conditioned over the observed data and current estimates of the parameters. Previous work on EM algorithm\cite{houdouin2023regularizedemalgorithm}.
    \item 
\end{itemize}

    
\end{document}